% ========================================
% CAPÍTULO 6: APLICABILIDAD Y USO RESPONSABLE  (criterio 6)
% ========================================

\chapter{APLICABILIDAD Y USO RESPONSABLE}
\label{ch:aplicabilidad}

\section{Contexto de uso}

Los usuarios previstos son (1) analistas y estudiantes de ciencia de datos que reciben pipelines de un asistente de IA y (2) equipos que quieren una revisión automática antes de integrar código a un repositorio. La decisión que informa el auditor es: \emph{aceptar, corregir o derivar a revisión humana} un pipeline.

\section{Propuesta de implementación}

Se ofrecen dos vías equivalentes, que comparten el mismo auditor.

\begin{itemize}
  \item \textbf{Skills para agentes} (\arch{05_skills/}). \texttt{audit-ml-leakage} audita y repara scripts existentes; \texttt{create-\allowbreak leakage-\allowbreak free-\allowbreak pipeline} impone reglas obligatorias al crear código (todo preprocesamiento dentro de un \texttt{Pipeline}, prueba final apartada, semillas fijas) y exige autoverificación con el auditor antes de entregar. Cada Skill es una carpeta autocontenida que se copia al directorio de Skills del agente.
  \item \textbf{Plataforma web} (\arch{06_web/}, Streamlit). El usuario pega o sube un \texttt{.py}, ve las líneas con fuga, descarga el código corregido y consulta los resultados del experimento. Está pensada para quienes no usan agentes y \textbf{no ejecuta el código}: solo lo lee.
  \item \textbf{Integración continua.} La herramienta de línea de comandos devuelve código de salida 0 (limpio) o 1 (fuga), por lo que puede bloquear una integración.
\end{itemize}

\emph{Requisitos:} Python (probado con la versión 3.13) y scikit-learn. \emph{Condiciones de uso:} scripts tabulares con scikit-learn; los notebooks se exportan a \texttt{.py}. \emph{Monitoreo:} registrar los falsos negativos que reporten los usuarios y usarlos para ampliar las reglas. \emph{Verificación hecha:} 37 pruebas automáticas, incluida la comprobación de que la plataforma web arranca sin errores. \textbf{Aún no se ha hecho una prueba de uso con personas}; el producto es un prototipo y esa prueba queda para Capstone~II.

\section{Riesgos éticos y medidas}

\begin{itemize}
  \item \textbf{Falsa seguridad.} Un resultado «0 fugas detectadas» no significa «sin fuga»: no se ven fugas temporales, de dominio ni por duplicados, ni código dinámico. La plataforma y las Skills lo declaran en pantalla y en sus instrucciones.
  \item \textbf{Datos.} No se usan datos personales. \texttt{breast\_cancer} y \texttt{diabetes} son conjuntos clínicos públicos y anonimizados que se emplean solo con fines metodológicos, sin ninguna intención de uso clínico.
  \item \textbf{Seguridad.} El banco ejecuta solo código generado por el propio proyecto; el código de LLM reales se ejecutará únicamente en entornos aislados.
  \item \textbf{Sesgo.} No se evalúa equidad porque no es el objeto del estudio; se declara como límite.
  \item \textbf{Términos de servicio.} Se revisarán los de las APIs de LLM antes de publicar salidas generadas.
  \item \textbf{Uso de IA.} \textbf{[COMPLETAR según la política del programa: indicar qué partes se redactaron o programaron con asistencia de IA y que el autor las revisó y puede explicarlas.]}
\end{itemize}
